\scanpage{27}
Soit $f$ comme dans Cor. 2, \inserted{mais on suppose $f$ un \unsure{épim|épi|épim.} effectif.} On veut prouver que $f$ est un morphisme de descente \uline{effective}. Soit donc $A'$ schéma abélien sur $S'$, avec donnée de descente relativement à $S'\to S$. à prouver que cette donnée est effective.

a) Réduction au cas $S$ local \inserted{$=\operatorname{Spec} R$} : \uline{triviale}, par passage à la limite.

\deleted{On y \unsure{descend|se ramène} aisément, en \unsure{remarquant|disant} qu'il \unsure{suffit} pour \unsure{suite|voir} que la donnée est \unsure{effectivement|effective dans le} \illegible{4}.}

b) Réduction au cas $S$ local complet à corps résiduel alg. clos. Grâce au critère de descente plate de schémas abéliens de Raynaud. Mais il faut savoir que

\diagram{\begin{tikzcd}[ampersand replacement=\&,column sep=small,row sep=small]
S\arrow[r,dash]\arrow[d,dash] \& S'\arrow[r,dash,shift left=0.325]\arrow[r,dash,shift right=0.325]\arrow[d,dash] \& S''\arrow[r,dash,shift left=0.65]\arrow[r,dash]\arrow[r,dash,shift right=0.65]\arrow[d,dash] \& S'''\arrow[d,dash] \& \text{1-descente} \\
S_{1}\arrow[r,dash]\arrow[d,dash,shift left=0.325]\arrow[d,dash,shift right=0.325] \& S'_{1}\arrow[r,dash,shift left=0.325]\arrow[r,dash,shift right=0.325]\arrow[d,dash,shift left=0.325]\arrow[d,dash,shift right=0.325] \& S''_{1}\arrow[r,dash,shift left=0.65]\arrow[r,dash]\arrow[r,dash,shift right=0.65]\arrow[d,dash,shift left=0.325]\arrow[d,dash,shift right=0.325] \& S'''_{1}\arrow[d,dash,shift left=0.325]\arrow[d,dash,shift right=0.325] \& \text{1-descente} \\
S_{2}\arrow[r,dash,"f_2"']\arrow[d,dash,shift left=0.65]\arrow[d,dash]\arrow[d,dash,shift right=0.65] \& S'_{2}\arrow[r,dash,shift left=0.325]\arrow[r,dash,shift right=0.325]\arrow[d,dash,shift left=0.65]\arrow[d,dash]\arrow[d,dash,shift right=0.65] \& S''_{2}\arrow[r,dash,shift left=0.65]\arrow[r,dash]\arrow[r,dash,shift right=0.65]\arrow[d,dash,shift left=0.65]\arrow[d,dash]\arrow[d,dash,shift right=0.65] \& S'''_{2}\arrow[d,dash,shift left=0.65]\arrow[d,dash]\arrow[d,dash,shift right=0.65] \& \text{1-descente} \\
S_{3}\arrow[r,dash] \& S'_{3}\arrow[r,dash,shift left=0.325]\arrow[r,dash,shift right=0.325] \& S''_{3}\arrow[r,dash,shift left=0.65]\arrow[r,dash]\arrow[r,dash,shift right=0.65] \& S'''_{3} \& \text{0-descente}
\end{tikzcd}}

$S'_2\to S_2$ est \inserted{morphisme} de descente pour la catégorie des schémas abéliens relatifs, et on ne peut directement appliquer le \deleted{Lemme 1, car} cor. 2, car $S_2$ n'est plus loc. noeth ! Mais on s'en tire en remarquant qu'il est vrai \deleted{que l'on} ici que (avec les notations de la dém. du Lemme 1) \deleted{$A$ est radiciel sur \unsure{$\operatorname{Ker}h_{2}$|$\operatorname{Ker}h_{2*}$}} $\operatorname{Ker}(f_{2*}(\mathcal O_{S'_2})\rightrightarrows h_{2*}(\mathcal O_{S''_2}))$ est radiciel sur $\mathcal O_{S_2}$, car tout provient par chgt de base plat de la situation $\mathcal O_S\to\operatorname{Ker}(f_*(\mathcal O_{S'})\rightrightarrows h_*(\mathcal O_{S''}))$. \unsure{un|en}
% slanted
\margin{dernière remarque après Lemme 1, donne aussi le cas de 0-descente}

c) \deleted{\unsure{Cas} $S$ \unsure{algébrique|alg. clos}} Cas \deleted{en $\dim S=0$} \inserted{\unsure{$S$ local de $\dim S=0$|$S$ local de $\dim{}\leq0$} \deleted{(et \unsure{$S$ local corps résiduel|corps résiduel} alg. clos)}}

\deleted{On se ramène au cas où $S'$ est fini sur $S$ (quitte à prendre un sous-schéma fini \deleted{de $S'$} assez grand de $S'$), puis \uline{par factorisation}, \uline{\unsure{successives|successivement}}, au cas où \deleted{\unsure{long}} $\lg_A A'/A=1$. Alors le critère de descente non plate nous ramène à trouver des ouverts affines recouvrant $A'$, stables par \deleted{des} donnée de descente. Or pour ça, on est ramené au cas où $A$ est le spectre d'un corps, où c'est trivial.}
% slanted
\margin{\deleted{expliciter des \unsure{points|point} de vue ``\unsure{General|Generalised} \unsure{non-sense|non sense}''}}

\deleted{d) Construction d'un schéma \inserted{abélien} formel sur $\operatorname{Spf} R$, qui descend le schéma formel défini par $A'$ (complété formel le long de $A'_s$ \unsure{\dots}). \deleted{Immédiat par c)} \inserted{Facile}.}

% archive sheet number: 26

